%----------------------------------------------------------------------
\section{Background}\label{sec:setting}
%----------------------------------------------------------------------
\label{sec:fv}

% Lean~4~\cite{demoura2021} is an interactive theorem prover and programming
% language built on a theory known as the Calculus of Inductive Constructions.
% Essentially, this allows it to represent mathematical structures and arguments
% in type theory (as opposed to the standard Set Theory) and verify the
% correctness of statements. Once Lean has verified that a statement is correct,
% to be fully confident in its output you must only believe two things:
%
% \begin{enumerate}[leftmargin=1.5em]
%     \item That the formal statement in Lean is equivalent to the informal
% counterpart
%     \item That Lean's small, independently verified proof kernel is correct.
% \end{enumerate}
% Thus, verifying mathematical statements in Lean gives you significantly more
% confidence in the truth of the statement. So long as the statements match,
% there is no need for a human to check the arguments used to prove the
% statements.
%
% One of the main reasons for using Lean is its robust and wide-ranging Mathlib
% library, which provides nearly 2 million lines of code across many branches of
% mathematics. Mathlib is able to define an increasing amount of research
% mathematics, and underpins all major Lean formalization projects like the
% Prime Number Theorem \cite{Kontorovich_Prime_Number_Theorem_2024} and Fermat's
% Last Theorem projects \cite{FLT_Lean}. For our project, Mathlib provides a
% number of indispensable algebraic objects that support the theory of Quantum
% Error Collection, and any large Lean project must plan to have some of its
% work eventually adopted into Mathlib for use by others.
%
% To provide an intuition on interactive theorem proving, we will sketch one of
% our lemmas and provide a side-by-side comparison with the Lean formalization.
% In general, Lean formalizations follow the same structure as pen-and-paper
% arguments, except all minor details must be fleshed out.
% For example, while formalizing code distances,
% since we cannot take a minimum over the empty set,
% we must explicitly handle an edge case for when there are no undetectable
% Pauli errors.
% (That is, when the only valid codeword is $0^n$.)
% Concretely, we define the \emph{distance} of a quantum error correcting code
% as follows.
%

This section gives some background Lean and formal verification for
\Cref{sec:QECBackground,sec:smt}. We assume familiarity with quantum error
correction and standard linear algebra.
For more details please refer to \cite{avigad2020mathematics}.

\paragraph{Lean 4 and Mathlib.} Lean~4~\cite{demoura2021} is an interactive
theorem prover (ITP) and a dependently typed functional programming language
whose underlying logic is the Calculus of Inductive Constructions. Statements
and proofs are both first-class objects: a proof of a proposition $P$ is a value
whose type is $P$. Once Lean has accepted a proof, the only assumptions a reader
needs to discharge are (i)~that the Lean statement faithfully captures the
intended informal claim and (ii)~that Lean's small, independently audited kernel
is correct.

The standard library \emph{Mathlib}~\cite{mathlib2020} now contains close to two
million lines of formalized mathematics and underpins large efforts such as the
Prime Number Theorem~\cite{Kontorovich_Prime_Number_Theorem_2024} and Fermat's
Last Theorem~\cite{FLT_Lean} projects. We rely on Mathlib for matrices,
finite-dimensional vector spaces, group theory, unitary groups, and the
Kronecker product; conversely, our development is designed to be portable back
into Mathlib as quantum-information content matures there.

\paragraph{Formalizing Math with an ITP.} A Lean proof is a sequence of \emph{tactics} that incrementally transforms the goal, and any informal pen-and-paper argument can be translated by spelling out the steps the human reader would otherwise fill in. The friction sits in those omitted steps: edge cases that the prose dismisses with ``clearly,'' implicit coercions between mathematically equivalent representations, and definitions whose set-theoretic phrasing is awkward in type theory.

\begin{mybox}{\textbf{Example: Proving Distance in Lean}}

For a concrete Lean example we revisit a definitional edge case that arises in our own
development. The distance of a quantum error-correcting code is the minimum
Pauli weight of an undetectable error, but a minimum over the empty set is
undefined; we must pick a convention for codes whose only valid codeword is
$0^n$. We follow the standard ``$n+1$'' convention, which makes the distance
well-defined while reserving the literal value $n+1$ as a sentinel for the
trivial case.

\textbf{Definition:} Given an error-correcting code $\mathsf{C}$, let $S\subseteq \mathbb{P}_n$
    be the set of undetectable Pauli errors.
    We then define the distance by
    \begin{equation}
        d = \begin{cases}
            \min_{P\in S}(\mathsf{weight}(P)) & \textnormal{if } S\ne\emptyset \\
            n+1 & \textnormal{if } S = \emptyset
        \end{cases}
    \end{equation}


The sentinel value makes the following sanity check immediate, but only
\emph{after} the case split is made explicit.

\textbf{Lemma:} If the distance $d$ of a code over $n$ qubits satisfies $d \leq n$, then its
    set $S$ of undetectable errors is nonempty.
\begin{proof}
    Suppose $S = \emptyset$ for contradiction. By the definition of code distance, we
    have $d = n+1$, contradicting $d \leq n$.
\end{proof}

A corresponding Lean proof mirroring step-by-step the above informal argument would be written as:

\begin{lstlisting}
lemma BinSympMatrix.undetectable_set_nonempty {k n : ℕ}
    (B : BinSympMatrix k n) (hk : 0 < k)
    (h_dist : B.distance hk ≤ n) :
  B.undetectable_set.Nonempty := by
  by_contra h_empty
  rw [Finset.not_nonempty_iff_eq_empty] at h_empty
  unfold BinSympMatrix.distance at h_dist
  rw [h_empty] at h_dist
  simp at h_dist
\end{lstlisting}    

This proof may be read as follows: \lstinline{by_contra} introduces the
contradiction hypothesis; \lstinline{Finset.not_nonempty_iff_eq_empty} converts
``non-emptiness'' into a Finset-level equation that Mathlib has lemmas about;
\lstinline{unfold} and \lstinline{rw} apply the definition and the empty-set
hypothesis; \lstinline{simp} discharges the residual arithmetic against $d \le
n$. The mechanical overhead, which involves naming the empty-set lemma and finding the right lemmas to rewrite with, is representative of what formalizing a pen-and-paper QEC
argument feels like in practice. 


\end{mybox}


%
% Lean proofs are composed of a series of "tactics", which perform mathematical
% proof techniques like deductive logic identities, syntactic rewriting, and
% more complex simplification procedures. The theorem is stated as:
%
% \begin{lstlisting}
%     theorem dual_dual_eq {C : CodeSpace n} : (C.dualCode).dualCode = C
% \end{lstlisting}
%
% As an Interactive Theorem Prover, Lean provides a ``context window`` the
% displays the user's slate of current hypotheses, and the ``goal'' they are
% trying to prove in the proof environment. For this theorem, it looks like:
%
% \begin{lstlisting}
% n : ℕ
% C : CodeSpace n
% ⊢ C.dualCode.dualCode = C
% \end{lstlisting}
%
% We plan to prove it by way of the Mathlib lemma
% \lstinline{Submodule.eq_of_le_of_finrank_eq}, which has statement:
%
% \begin{lstlisting}
% Submodule.eq_of_le_of_finrank_eq.{u, v} {K : Type u} {V : Type v}
%   [DivisionRing K] [AddCommGroup V] [Module K V]
%   {S₁ S₂ : Submodule K V} [FiniteDimensional K ↥S₂] (hle : S₁ ≤ S₂)
%   (hd : Module.finrank K ↥S₁ = Module.finrank K ↥S₂) :
%   S₁ = S₂
% \end{lstlisting}
%
% This may be read as follows: beginning with \lstinline|{K : Type u}| are the
% arguments for the lemma, which requires types \lstinline{K} and \lstinline{U}
% with \lstinline{K} being a Division Ring, and \lstinline{V} being a
% commutative module over $K$. With types specified, submodules \lstinline{S₁, % S₂} are then
% defined, and \lstinline{S₂} is required to be finite dimensional.
% The two main hypotheses are that \lstinline{S₁} is a submodule of
% \lstinline{S₂}, and that \lstinline{S₁} and \lstinline{S₂} have equal rank.
% After the colon is the ``type'' of the lemma, or the proposition it
% ``produces'' a proof object for (by the Curry-Howard isomorphism, Lean proofs
% are functions in a functional programming language).
% \Ethan{I feel like either the reader already knows Curry-Howard or they don't.
% In the latter case, we have no chance of successfully conveying the idea.}
% Note that the lemma arguments are variously provided in square brackets, curly
% braces, and parentheses. The parentheses-closed arguments are explicitly
% provided to the function, while the remainder of the arguments are
% automatically synthesized.
%
% Running the following tactics:
%
% \begin{lstlisting}
% symm
% apply Submodule.eq_of_le_of_finrank_eq
% \end{lstlisting}
%
% First reverses the order of the equality in the goal, then ``applies'' the
% lemma, synthesizing its arguments with the goal and seeking explicit
% instantiations of the goal's arguments. The context window then looks like:
%
% \begin{lstlisting}
% 2 goals
% case hle
% n : ℕ
% C : CodeSpace n
% ⊢ C ≤ C.dualCode.dualCode
% case hd
% n : ℕ
% C : CodeSpace n
% ⊢ Module.finrank (ZMod 2) ↥C = Module.finrank (ZMod 2) ↥C.dualCode.dualCode
% \end{lstlisting}
%
% To dispatch the first goal, we need to prove that \lstinline{S ≤ % S.dualCode.dualCode}.
% \Ethan{I feel anything beyond this is too much.
% We should probably refer to some standard course/textbook instead.
% We aren't here to teach the readers how interactive theorem provers work.}
% An earlier definition specifies the dual as the orthogonal submodule with
% respect to the dot product, so we apply another Mathlib lemma about submodules
% defined as:
%
% \begin{lstlisting}
% Submodule.le_orthogonalBilin_orthogonalBilin.{u_1, u_2, u_5, u_6} {R : Type
% u_1} {R₁ : Type u_2} {M : Type u_5}
%   {M₁ : Type u_6} [CommRing R] [CommRing R₁] [AddCommGroup M₁] [Module R₁ M₁]
% [AddCommGroup M] [Module R M]
%   {I₁ I₂ : R₁ →+* R} {B : M₁ →ₛₗ[I₁] M₁ →ₛₗ[I₂] M} {N : Submodule R₁ M₁} (b :
% B.IsRefl) :
%   N ≤ (N.orthogonalBilin B).orthogonalBilin B
% \end{lstlisting}
%
% The other goal about equal rank is proved using an earlier lemma defining the
% rank of the dual code, \lstinline{dual_finrank_eq}.
%
% \begin{lstlisting}
% lemma dual_finrank_eq {C : CodeSpace n} : Module.finrank (ZMod 2) C.dualCode =
% n - (Module.finrank (ZMod 2) C)
% \end{lstlisting}
%
% Rewriting this twice, the goal is now:
%
% \begin{lstlisting}
% ⊢ Module.finrank (ZMod 2) ↥C = n - (n - Module.finrank (ZMod 2) ↥C)
% \end{lstlisting}
%
% Which is proven by simplification tactics after asserting that \lstinline{C}
% has rank at most $n$, as the subtraction is occurring over the natural numbers
% and would be false otherwise.
%
%

%This example demonstrates Lean's expressive power as well as the benefits
%gained from working with Mathlib.

%\subsection{Satisfiability Modulo Theories}

%\emph{Satisfiability modulo theories} (``SMT'') is a natural generalization of
%the Boolean satisfiability (``SAT'') problem.
%Instead of working with only Boolean variables, SMT allows formulas to include
%various data types and structures.
%Similar to SAT, SMT is NP-hard in general,
%thus SMT instances are usually treated using specialized solvers such as cvc5
%\cite{barbosa2022}, alt-ergo \cite{conchon2018alt}, and Z3 \cite{de2008z3}.
%The comparisons and inner workings of these solvers are out of the scope of our
%discussions.
%For the purpose of our project,
%we will treat solvers as black-boxes that run in (at worst) asymptotically
%exponential time.

%Most popular proof assistant software include various levels of solver support.
%In Lean, we use CaDiCaL \cite{biere2024cadical} for SAT instances, and Lean-SMT
%(based on cvc5) \cite{mohamed2025lean} for SMT instances.
%In these integrations,
%the solver should ideally not be included within the trusted computing base
%(TCB).
%Instead, the solver could return proofs of (un)satisfiability, which are then
%translated and verified using standard techniques such as Curry-Howard
%isomorphism. For cvc5, the standard solver contains a flag which causes it to
%return a "proof" of its reasoning in the Cooperating Proof Calculus (CPC),
%which may be verified by its checker Ethos and is compatible with the Eunoia
%framework. Unfortunately, this translation step from proof calculi to Lean, due
%to the overhead of being performed within Lean, becomes computationally
%prohibitive for larger and more complex SMT instances.
%As a result, frameworks such as EasyCrypt have chosen to absorb their solvers
%into their TCB, and to accept their external solvers' outputs without
%verification.

%On the proof engineering side, SMT solvers are often treated similarly to
%\emph{hammer} tactics.
%In other words, they are designed for goals that seem more tedious than
%complex.
%As a result, it is conventional to set a timeout of well under one minute per
%SMT invocation.
%As we discuss in \Cref{sec:smt} however,
%some of our project's goals require hours-long SMT runs.
